\section{Übung 10 – Funktion und Transformation von Zufallsvariablen}

In dieser Übung wird die Vorrechenaufgabe~10 gelöst. Es geht darum, aus der
bekannten Dichte einer Zufallsvariablen $X$ die Dichte einer transformierten
Zufallsvariablen $Y = g(X)$ zu bestimmen. Teil~A behandelt Funktionen einer
exponentialverteilten Zufallsvariablen (eine quadratische und eine lineare
Transformation), Teil~B die Standardisierung einer normalverteilten
Zufallsvariablen. Grundlage ist der Transformationssatz für streng monotone
Funktionen aus Abschnitt~\ref{sec:monotone-transformation}.

% ----------------------------------------------------------------
\subsection*{Teil A – Funktion einer exponentialverteilten Zufallsvariablen}
% ----------------------------------------------------------------

Die Zufallsvariable $X$ ist exponentialverteilt
(Abschnitt~\ref{sec:exponentialverteilung}) mit Parameter $\lambda > 0$:

$$
p_X(x) = \begin{cases} \lambda\, e^{-\lambda x} & , \quad x \geq 0 \\[2pt]
0 & , \quad \text{sonst.} \end{cases}
$$

Der Träger von $X$ ist also $[0, \infty)$. Diese Einschränkung ist für die
Wahl der Umkehrfunktion entscheidend.

% ---- A1 --------------------------------------------------------
\subsubsection*{A1 – Quadratische Transformation $Y = X^2$}

\begin{beispiel}[Funktion einer Zufallsvariablen -- $Y = X^2$]

\textbf{Gegeben:} Exponentialverteilte Zufallsvariable $X$ mit Dichte
$p_X(x) = \lambda e^{-\lambda x}$ für $x \geq 0$ und Parameter $\lambda > 0$.
Transformation $Y = g(X) = X^2$.

\textbf{Gesucht:} Die Wahrscheinlichkeitsdichte $p_Y(y)$ der transformierten
Zufallsvariablen $Y$.

\textbf{Formel} (Transformationssatz für streng monotone Funktionen,
Abschnitt~\ref{sec:monotone-transformation}):

$$
p_Y(y) = \frac{p_X\!\left(g^{-1}(y)\right)}{\left|\, g'\!\left(g^{-1}(y)\right)\right|}
$$

\textbf{Rechnung:}

Zwar ist $g(x) = x^2$ auf ganz $\mathbb{R}$ \emph{nicht} streng monoton, auf dem
Träger $x \geq 0$ von $X$ jedoch streng monoton steigend. Es kommt daher nur die
\emph{positive} Wurzel als Umkehrfunktion in Frage:

$$
x = g^{-1}(y) = \sqrt{y}, \qquad y \geq 0.
$$

Die Ableitung von $g(x) = x^2$ ist $g'(x) = 2x$, an der Stelle $x = \sqrt{y}$
also

$$
\left|\, g'\!\left(\sqrt{y}\right)\right| = 2\sqrt{y}.
$$

Einsetzen in den Transformationssatz liefert für $y \geq 0$:

$$
p_Y(y) = \frac{p_X\!\left(\sqrt{y}\right)}{2\sqrt{y}}
       = \frac{\lambda\, e^{-\lambda \sqrt{y}}}{2\sqrt{y}}.
$$

Für $y < 0$ existiert kein reelles Urbild, also $p_Y(y) = 0$.

\textbf{Lösung:}

$$
\boxed{\;p_Y(y) = \begin{cases}
\dfrac{\lambda}{2\sqrt{y}}\, e^{-\lambda \sqrt{y}} & , \quad y > 0 \\[8pt]
0 & , \quad \text{sonst.}
\end{cases}\;}
$$

An der Stelle $y \to 0^+$ besitzt die Dichte wegen des Faktors
$1/(2\sqrt{y})$ eine (integrierbare) Polstelle -- die Wahrscheinlichkeitsmasse
wird durch das Quadrieren nahe null stark „gestaucht``. Zur Kontrolle gilt
weiterhin $\int_0^\infty p_Y(y)\,\mathrm{d}y = 1$ (Substitution
$y = x^2$ führt auf das Integral der Ausgangsdichte).

\end{beispiel}

% ---- A2 --------------------------------------------------------
\subsubsection*{A2 – Lineare Transformation $Y = aX + b$}

\begin{beispiel}[Funktion einer Zufallsvariablen -- $Y = aX + b$]

\textbf{Gegeben:} Exponentialverteilte Zufallsvariable $X$ mit Dichte
$p_X(x) = \lambda e^{-\lambda x}$ für $x \geq 0$. Lineare Transformation
$Y = g(X) = aX + b$ mit Konstanten $a \neq 0$ und $b \in \mathbb{R}$.

\textbf{Gesucht:} Die Wahrscheinlichkeitsdichte $p_Y(y)$ der transformierten
Zufallsvariablen $Y$.

\textbf{Formel} (Transformationssatz für streng monotone Funktionen,
Abschnitt~\ref{sec:monotone-transformation}):

$$
p_Y(y) = \frac{p_X\!\left(g^{-1}(y)\right)}{\left|\, g'\!\left(g^{-1}(y)\right)\right|}
$$

\textbf{Rechnung:}

Die lineare Funktion $g(x) = ax + b$ ist wegen $a \neq 0$ streng monoton und
eindeutig umkehrbar:

$$
x = g^{-1}(y) = \frac{y - b}{a}, \qquad g'(x) = a, \qquad |g'(x)| = |a|.
$$

Einsetzen ergibt zunächst formal

$$
p_Y(y) = \frac{1}{|a|}\, p_X\!\left(\frac{y - b}{a}\right).
$$

Die Ausgangsdichte $p_X$ ist nur für ein nichtnegatives Argument von null
verschieden, also für $\dfrac{y - b}{a} \geq 0$. Für den typischen Fall
$a > 0$ bedeutet das $y \geq b$, und man erhält

$$
p_Y(y) = \frac{\lambda}{a}\,
         \exp\!\left(-\lambda\,\frac{y - b}{a}\right), \qquad y \geq b.
$$

\textbf{Lösung} (Fall $a > 0$):

$$
\boxed{\;p_Y(y) = \begin{cases}
\dfrac{\lambda}{a}\, \exp\!\left(-\dfrac{\lambda}{a}\,(y - b)\right)
  & , \quad y \geq b \\[8pt]
0 & , \quad y < b.
\end{cases}\;}
$$

$Y$ ist damit wieder \emph{exponentialverteilt}, allerdings um $b$ nach rechts
verschoben und mit dem neuen Parameter $\lambda' = \lambda/a$. Der Faktor $a$
skaliert die Zufallsvariable und verkleinert entsprechend den Ratenparameter
(im Fall $a < 0$ kehrt sich der Träger zu $y \leq b$ um; die kompakte
Bedingung lautet allgemein $\tfrac{y-b}{a} \geq 0$, und der Vorfaktor ist
$\lambda/|a|$).

\end{beispiel}

% ----------------------------------------------------------------
\subsection*{Teil B – Standardisierung einer normalverteilten Zufallsvariablen}
% ----------------------------------------------------------------

% ---- B1 --------------------------------------------------------
\subsubsection*{B1 – Transformation auf die Standardnormalverteilung}

\begin{beispiel}[Standardisierung $X \sim \mathcal{N}(\mu_X, \sigma_X^2) \to Y \sim \mathcal{N}(0,1)$]

\textbf{Gegeben:} Normalverteilte Zufallsvariable $X$
(Abschnitt~\ref{sec:normalverteilung}) mit Mittelwert $\mu_X$ und Varianz
$\sigma_X^2 > 0$:

$$
p_X(x) = \frac{1}{\sqrt{2\pi}\,\sigma_X}\,
         \exp\!\left(-\frac{(x - \mu_X)^2}{2\sigma_X^2}\right).
$$

\textbf{Gesucht:} Eine lineare Transformation $Y = g(X)$, so dass $Y$
standardnormalverteilt ist ($\mu_Y = 0$, $\sigma_Y^2 = 1$), sowie der Nachweis
über die transformierte Dichte $p_Y(y)$.

\textbf{Formel} (Standardisierung sowie Transformationssatz,
Abschnitt~\ref{sec:monotone-transformation}):

$$
Y = \frac{X - \mu_X}{\sigma_X}, \qquad
p_Y(y) = \frac{p_X\!\left(g^{-1}(y)\right)}{\left|\, g'\!\left(g^{-1}(y)\right)\right|}
$$

\textbf{Rechnung:}

\emph{Ansatz.} Die Standardisierung subtrahiert den Mittelwert und normiert auf
die Standardabweichung:

$$
Y = g(X) = \frac{X - \mu_X}{\sigma_X}.
$$

Dass diese Wahl die geforderten Momente erzeugt, folgt direkt aus der
Linearität von Erwartungswert (Abschnitt~\ref{sec:erwartungswert}) und Varianz:

$$
\mu_Y = \mathrm{E}\{Y\} = \frac{\mathrm{E}\{X\} - \mu_X}{\sigma_X} = 0,
\qquad
\sigma_Y^2 = \mathrm{var}\{Y\} = \frac{\mathrm{var}\{X\}}{\sigma_X^2}
           = \frac{\sigma_X^2}{\sigma_X^2} = 1.
$$

\emph{Dichte.} Die Funktion $g(x) = (x - \mu_X)/\sigma_X$ ist streng monoton
steigend mit Umkehrfunktion und Ableitung

$$
x = g^{-1}(y) = \sigma_X\, y + \mu_X, \qquad
g'(x) = \frac{1}{\sigma_X}, \qquad |g'(x)| = \frac{1}{\sigma_X}.
$$

Einsetzen in den Transformationssatz:

$$
p_Y(y) = \frac{p_X\!\left(\sigma_X\, y + \mu_X\right)}{1/\sigma_X}
       = \sigma_X\, p_X\!\left(\sigma_X\, y + \mu_X\right).
$$

Nun wird $x = \sigma_X y + \mu_X$ in die Normaldichte eingesetzt; im Exponenten
gilt $(x - \mu_X)^2 = (\sigma_X y)^2 = \sigma_X^2\, y^2$:

$$
p_Y(y) = \sigma_X \cdot \frac{1}{\sqrt{2\pi}\,\sigma_X}\,
         \exp\!\left(-\frac{\sigma_X^2\, y^2}{2\sigma_X^2}\right)
       = \frac{1}{\sqrt{2\pi}}\, \exp\!\left(-\frac{y^2}{2}\right).
$$

\textbf{Lösung:}

$$
\boxed{\;Y = \frac{X - \mu_X}{\sigma_X}
\quad\Longrightarrow\quad
p_Y(y) = \frac{1}{\sqrt{2\pi}}\, e^{-y^2/2}
= \mathcal{N}(y;\, 0,\, 1)\;}
$$

$Y$ ist also standardnormalverteilt. Die Standardabweichung $\sigma_X$ aus dem
Betrag der Ableitung kürzt sich exakt gegen den Vorfaktor der Ausgangsdichte,
und die Verschiebung um $\mu_X$ verschwindet im Exponenten. Diese
Standardisierung ist die Grundlage für die Nutzung tabellierter Werte der
Standardnormalverteilung $\Phi(y)$ für beliebige normalverteilte Größen.

\end{beispiel}

% ----------------------------------------------------------------
\subsection*{Aufgabe 10.1 – Transformation auf eine Standard-Gleichverteilung}
% ----------------------------------------------------------------

\begin{beispiel}[Gleichverteilung auf einem Intervall zur Standard-Gleichverteilung]

\textbf{Gegeben:} Eine auf dem Intervall $[a, b]$ (mit $a < b$) gleichverteilte
Zufallsvariable $X$ (Abschnitt~\ref{sec:gleichverteilung}):

$$
p_X(x) = \begin{cases} \dfrac{1}{b - a} & , \quad a \leq x \leq b \\[6pt]
0 & , \quad \text{sonst.} \end{cases}
$$

\textbf{Gesucht:} Eine Transformation $Y = g(X)$, so dass $Y$ auf dem Intervall
$\left[-\tfrac{1}{2}, \tfrac{1}{2}\right]$ gleichverteilt ist, sowie der
Nachweis über die transformierte Dichte $p_Y(y)$.

\textbf{Formel} (lineare Transformation $Y = cX + d$ und Transformationssatz,
Abschnitt~\ref{sec:monotone-transformation}):

$$
p_Y(y) = \frac{p_X\!\left(g^{-1}(y)\right)}{\left|\, g'\!\left(g^{-1}(y)\right)\right|}
$$

\textbf{Rechnung:}

\emph{Ansatz.} Da beide Intervalle die Endpunkte einer Gleichverteilung
festlegen, genügt eine \emph{lineare} und streng monoton steigende
Transformation $Y = cX + d$, welche die Intervallgrenzen aufeinander abbildet:

$$
g(a) = -\tfrac{1}{2}, \qquad g(b) = +\tfrac{1}{2}.
$$

Die Steigung ergibt sich aus dem Verhältnis der Intervalllängen:

$$
c = \frac{\tfrac{1}{2} - \left(-\tfrac{1}{2}\right)}{b - a} = \frac{1}{b - a}.
$$

Der Versatz $d$ folgt aus $g(a) = -\tfrac{1}{2}$:

$$
d = -\tfrac{1}{2} - \frac{a}{b - a} = -\frac{a + b}{2\,(b - a)}.
$$

Zusammengefasst lässt sich die Transformation kompakt als Zentrierung um den
Mittelwert $\mu_X = \tfrac{a+b}{2}$ und Normierung auf die Intervalllänge
schreiben:

$$
Y = g(X) = \frac{X - \dfrac{a + b}{2}}{b - a}.
$$

\emph{Kontrolle der Grenzen:}

$$
g(a) = \frac{a - \tfrac{a+b}{2}}{b - a} = \frac{-\tfrac{b-a}{2}}{b-a} = -\tfrac{1}{2},
\qquad
g(b) = \frac{b - \tfrac{a+b}{2}}{b - a} = \frac{\tfrac{b-a}{2}}{b-a} = +\tfrac{1}{2}.
\quad\checkmark
$$

\emph{Dichte.} Die Funktion $g$ ist streng monoton steigend mit

$$
x = g^{-1}(y) = (b - a)\, y + \frac{a + b}{2}, \qquad
g'(x) = \frac{1}{b - a}, \qquad |g'(x)| = \frac{1}{b - a}.
$$

Einsetzen in den Transformationssatz für $y \in \left[-\tfrac12, \tfrac12\right]$
(dort ist $g^{-1}(y) \in [a, b]$, also $p_X = \tfrac{1}{b-a}$):

$$
p_Y(y) = \frac{p_X\!\left(g^{-1}(y)\right)}{|g'|}
       = \frac{\dfrac{1}{b - a}}{\dfrac{1}{b - a}}
       = 1.
$$

\textbf{Lösung:}

$$
\boxed{\;Y = \frac{X - \dfrac{a + b}{2}}{b - a}
\quad\Longrightarrow\quad
p_Y(y) = \begin{cases} 1 & , \quad -\tfrac{1}{2} \leq y \leq \tfrac{1}{2} \\[4pt]
0 & , \quad \text{sonst.} \end{cases}\;}
$$

$Y$ ist also auf $\left[-\tfrac{1}{2}, \tfrac{1}{2}\right]$ gleichverteilt; die
Dichte hat dort den Wert $1$, da das Zielintervall die Länge $1$ besitzt. Die
Transformation entspricht einer Standardisierung der Gleichverteilung: Sie
verschiebt den Mittelwert nach $0$ und skaliert die Spannweite auf $1$
(analog zur Standardisierung der Normalverteilung in Teil~B, dort jedoch mit
$\sigma_X$ statt der Intervalllänge $b - a$).

\end{beispiel}
